\section*{Chapter \ref{ch:qm:games-auctions-and-information} --- Games, Auctions and Information}

\begin{solution}{exo:qm:games-auctions-and-information:1}
$b(v) = \frac{n-1}nv = 0.9 \times 0.7 = 0.63$; expected revenue $\frac{n-1}{n+1} = \frac9{11} = 0.818$.
\end{solution}

\begin{solution}{exo:qm:games-auctions-and-information:2}
Each player mixes $(\frac12, \frac12)$: any other mix gives the opponent a pure best response that wins on average. The value is 0.
\end{solution}

\begin{solution}{exo:qm:games-auctions-and-information:3}
$H = \log_26 = 2.585$ bits. Parity is a function of the die, so $I = H(\text{parity}) - H(\text{parity} \mid \text{die}) = 1 - 0 = 1$ bit.
\end{solution}

\begin{solution}{exo:qm:games-auctions-and-information:4}
$\int_0^vF^{n-1} = \int_0^vx^{2n-2}dx = v^{2n-1}/(2n-1)$, and dividing by $F(v)^{n-1} = v^{2n-2}$ gives $v/(2n-1)$, so $b(v) = v\frac{2n-2}{2n-1}$: with five bidders $\frac89v$, less shading than with uniform values
because rivals' values are more often high. The general formula, evaluated numerically, gives 0.2667, 0.5333 and 0.8000 at 0.3, 0.6 and 0.9.
\end{solution}

\begin{solution}{exo:qm:games-auctions-and-information:5}
$r - (1 - r) = 0.2$, so $r^* = 0.6$.
\end{solution}

\begin{solution}{exo:qm:games-auctions-and-information:6}
With $k = 1$, the bid is $E[Y \mid Y < v]$ with $Y$ the highest of the other $n - 1$ values, whose distribution function is $F^{n-1}$; $E[Y \mid Y < v] = v - \int_0^vF^{n-1}/F(v)^{n-1}$ by integration by parts, the
first-price bid.
\end{solution}

\begin{solution}{exo:qm:games-auctions-and-information:7}
\texttt{equivalence()}: 0.6665 (first), 0.6665 (second), 0.6715 (English, ticks of 0.01), 0.6614 (Dutch, ticks of 0.01), against $\frac23$; with the reserve 0.5 the \omterm{def:qm:games-auctions-and-information:auctions}{second-price auction} gives 0.6720 against the closed
form 0.6719, selling 96.9\% of the time.
\end{solution}

\begin{solution}{exo:qm:games-auctions-and-information:8}
Winning every competition at fair value means the firm's valuations were the highest each time, which is exactly when a noisy valuation overstates the truth: the winner's curse. A win rate near 100\%
against four rivals says the bids are too high, not that the model is good; the right check is the realised profit on the trades won.
\end{solution}

\begin{solution}{pb:qm:games-auctions-and-information:1}
\textbf{1.} $\frac45$ of its value: the equilibrium trades the probability of winning against the margin, and recovers the expected gap to the next value.
\textbf{2.} The bid only decides whether one wins, not what one pays; deviating wins only at a loss or loses only a profit.
\textbf{3.} Expected profit peaks at a bid of 0.64, the equilibrium bid: the strategy is a best response.
\textbf{4.} $\frac23$ in every format; one million auctions give 0.6665 (first price) and 0.6665 (second price).
\textbf{5.} The English clock stops at the first tick above the second value (half a tick high on average), the Dutch clock at the first tick below the bid (half a tick low): 0.6715 and 0.6614.
\textbf{6.} $r^* = 0.5$: from $\frac13$ to $\frac5{12}$ (25\%) with two bidders, unsold a quarter of the time; from 0.6667 to 0.6719 (0.8\%) with five, unsold 3.1\% of the time.
\textbf{7.} $b(v) = v(3 - 2.4v)/(4 - 3v)$: 0.36 at 0.5, 0.54 at 0.8, 0.60 at 1.
\textbf{8.} 1 in expectation for both; simulation gives 0.9994 (uniform) and 0.9999 (pay-as-bid).
\textbf{9.} The uniform format spreads bids about twice as widely (winning-bid gap 0.167 against 0.087) and its revenue is more than twice as volatile (0.378 against 0.160): more aggressive, more dispersed bidding
and more auction-to-auction volatility, as the study reported.
\textbf{10.} With unit demand, symmetric risk-neutral bidders and independent values, revenue equivalence makes the formats equal on average; the answer depends on multi-unit demand, risk aversion, asymmetry and
information, which is why the Treasury ran an experiment.
\textbf{11.} Bidding one's signal loses 0.067 per win (the expected overstatement of the highest of five signals), in 97\% of wins; shading by $0.1 \times \frac46$ brings the profit to zero.
\textbf{12.} 1.846 bits; $-0.165$ bits a race (the take exceeds the bettor's knowledge).
\textbf{13.} 0.363 bits, turning the doubling rate to $+0.198$; simulation gives a gain of $0.363 \pm 0.002$.
\textbf{14.} They multiply both doubling rates by the same factor inside the logarithm and cancel in the difference.
\textbf{15.} Nothing: with four horses a tip right with probability $\frac14$, and uniformly wrong otherwise, is independent of the winner.
\textbf{16.} Bidders' risk aversion or multi-unit demand reduction under pay-as-bid, a wider participation that the single price attracts, and a smaller winner's curse when every winner pays the same price.
\textbf{17.} Whenever others compete on the same information: shade by the expected overstatement of the winning estimate, more with more competitors and noisier valuations.
\textbf{18.} Less than the growth it adds: at most its \omterm{def:qm:games-auctions-and-information:entropy}{mutual information} with the outcome, in doubling rate, net of what the tip costs per race.
\textbf{19.} Named result: \emph{the Treasury's switch}: in the symmetric private-value model both formats raise the same expected revenue (1 for two units and five dealers; $\frac23$ in the single-unit case), each dealer
shades its first-price bid by one fifth, and the optimal reserve of 0.5 adds 0.8\% with five dealers (25\% with two).
\textbf{20.} Under its conditions the format does not matter for expected revenue; what matters is who participates, what they know and the reserve.
\end{solution}

\begin{solution}{iq:qm:games-auctions-and-information:1}
Bid $v/2$: against a rival bidding $w/2$, bidding $b$ wins with probability $2b$ and earns $(v - b)2b$, maximised at $b = v/2$.
\iqlookfor{The trade-off and the answer $v/2$.}
\end{solution}

\begin{solution}{iq:qm:games-auctions-and-information:2}
In a common-value setting the winner is the bidder with the most optimistic estimate, so conditional on winning the estimate is biased upwards. Bid as if your estimate were the highest: subtract the
expected overstatement of the maximum of the competitors' estimates, which grows with their number and noise.
\iqlookfor{Conditioning on winning.}
\end{solution}

\begin{solution}{iq:qm:games-auctions-and-information:3}
Any auction that allocates to the highest value and gives zero surplus to the lowest type yields the same expected revenue, with independent private values and risk-neutral symmetric bidders. It fails with
risk aversion (first price raises more), affiliated or common values (English raises more), asymmetric bidders, budget constraints and multi-unit demand.
\iqlookfor{Conditions and two failures.}
\end{solution}

\begin{solution}{iq:qm:games-auctions-and-information:4}
Its \omterm{def:qm:games-auctions-and-information:entropy}{mutual information} with the outcome, in doubling rate: the optimal log growth rises by exactly $I(X; Y)$ per bet, independent of the odds.
\iqlookfor{\omterm{def:qm:games-auctions-and-information:entropy}{Mutual information}.}
\end{solution}

\begin{solution}{iq:qm:games-auctions-and-information:5}
No saddle point (maximin $-1$, minimax 1). Row mixes $x = \frac37$ on the first row, column $y = \frac27$ on the first column; the value is $\frac17$.
\iqlookfor{Indifference conditions.}
\end{solution}

\begin{solution}{iq:qm:games-auctions-and-information:6}
In a private-value first-price setting more competitors mean less shading (bid closer to value); with common values (most of the portfolio's value is the same for all dealers) more competitors mean a worse winner's
curse and more shading. A portfolio request-for-quote has a large common component, so the second effect usually dominates.
\iqlookfor{The two opposite effects.}
\end{solution}
